\section{假设、数据处理与符号}
\subsection{题设口径与补充假设}
所有需求、时限和设备参数取自官方五份工作簿，地形取官方 30 m DEM。坐标为经纬度，水平距离采用 WGS84 椭球测地线距离\cite{karney}；航迹与 DEM 穿越按代码一致的经纬度端点线性插值计算。该短距离近似保持各模块口径一致，但不代表高精度航空导航轨迹。

依据题设，货箱不可拆分，每箱恰好交付一次；每个运输架次由 O01 出发并最终返回 O01。交付后对应货箱质量从剩余载荷中扣除。下降不单计附加能量，运输交接不另加未经题设给定的悬停功率。中继建链阶段按悬停与通信功率计入能耗。所有能源资源初始满电，再次投入前充至满电，不同资源可并行充电。

赛题未提供风速、风向或阵风参数，故采用给定计划速度，不另建风修正模型。DEM 视为像元内分段常数；线路恰好经过边界或角点时，纳入全部相接像元，以较高地形进行保守检查。中继窗口在已发现的直连不足区间两侧各扩展 5 s，这是明确的时序保护假设，不是连续通信证明。

\subsection{数据检查与符号系统}
读取后先检查编号唯一性、货箱归属、资源类型和数值范围，再检查必需航段和通信射线是否超出 DEM 或经过无效像元。正式计算使用完整精度；论文的三位能耗小数、两位载荷小数和秒级时间展示仅用于阅读。关键运输参数见表\ref{tab:parameters}，完整参数保留于原始工作簿。
\papertable{drone_parameters}{三种运输机型的主要参数；电池可用能量与额定载荷分别约束航程和装载}{parameters}

集合 $\mathcal S$ 表示服务区，$\mathcal B$ 表示货箱，$\mathcal G$ 表示机型，$\mathcal P$ 和 $\mathcal R$ 分别表示运输及中继架次。$i,j$ 为节点，$b$ 为货箱，$g$ 为机型，$p,r$ 为架次索引。表\ref{tab:symbols}区分参数、决策与派生状态。
\begin{table}[htbp]\centering\caption{核心符号及单位}\label{tab:symbols}
\begin{tabularx}{\linewidth}{lXl}\toprule
符号 & 含义与类别 & 单位\\\midrule
$d_{ij},z_i,H_{ij}$ & 水平距离、节点作业海拔、计划巡航海拔 & m\\
$m_b,v_b,w_b$ & 货箱质量、体积、应急优先系数（参数） & kg，m$^3$，1\\
$Q_g,V_g,E_g,\rho_g$ & 额定载荷、体积、电池可用能量、返航储备 & kg，m$^3$，kWh，1\\
$q_{p,ij},s_p$ & 航段剩余载荷（状态）、架次开始时刻（决策） & kg，s\\
$x_{pb},u_{pd},v_{pk}$ & 货箱、实体无人机、电池的分配变量 & 0--1\\
$C_b,\ell_b,T$ & 交付完成时刻、延误、任务完工历时（派生量） & s\\
$a_r,b_r,\bm h_r$ & 中继建链完成、服务结束、三维悬停位置 & s，s，经纬度/m\\
$M_{ij}(t)$ & 双向链路预算减去实际传播损耗（余量） & dB\\\bottomrule
\end{tabularx}\end{table}
\FloatBarrier
